% Options for packages loaded elsewhere
\PassOptionsToPackage{unicode}{hyperref}
\PassOptionsToPackage{hyphens}{url}
\documentclass[
  12pt,
]{article}
\usepackage{xcolor}
\usepackage{amsmath,amssymb}
\setcounter{secnumdepth}{-\maxdimen} % remove section numbering
\usepackage{iftex}
\ifPDFTeX
  \usepackage[T1]{fontenc}
  \usepackage[utf8]{inputenc}
  \usepackage{textcomp} % provide euro and other symbols
\else % if luatex or xetex
  \usepackage{unicode-math} % this also loads fontspec
  \defaultfontfeatures{Scale=MatchLowercase}
  \defaultfontfeatures[\rmfamily]{Ligatures=TeX,Scale=1}
\fi
\usepackage{lmodern}
\ifPDFTeX\else
  % xetex/luatex font selection
\fi
% Use upquote if available, for straight quotes in verbatim environments
\IfFileExists{upquote.sty}{\usepackage{upquote}}{}
\IfFileExists{microtype.sty}{% use microtype if available
  \usepackage[]{microtype}
  \UseMicrotypeSet[protrusion]{basicmath} % disable protrusion for tt fonts
}{}
\makeatletter
\@ifundefined{KOMAClassName}{% if non-KOMA class
  \IfFileExists{parskip.sty}{%
    \usepackage{parskip}
  }{% else
    \setlength{\parindent}{0pt}
    \setlength{\parskip}{6pt plus 2pt minus 1pt}}
}{% if KOMA class
  \KOMAoptions{parskip=half}}
\makeatother
\usepackage{longtable,booktabs,array}
\usepackage{calc} % for calculating minipage widths
% Correct order of tables after \paragraph or \subparagraph
\usepackage{etoolbox}
\makeatletter
\patchcmd\longtable{\par}{\if@noskipsec\mbox{}\fi\par}{}{}
\makeatother
% Allow footnotes in longtable head/foot
\IfFileExists{footnotehyper.sty}{\usepackage{footnotehyper}}{\usepackage{footnote}}
\makesavenoteenv{longtable}
\setlength{\emergencystretch}{3em} % prevent overfull lines
\providecommand{\tightlist}{%
  \setlength{\itemsep}{0pt}\setlength{\parskip}{0pt}}
\usepackage{bookmark}
\IfFileExists{xurl.sty}{\usepackage{xurl}}{} % add URL line breaks if available
\urlstyle{same}
\hypersetup{
  hidelinks,
  pdfcreator={LaTeX via pandoc}}

\author{SCX}
\date{2026}

\begin{document}

\section{Proposition 1: Regret Lower Bound for Global Expert
Selection}\label{proposition-1-regret-lower-bound-for-global-expert-selection}

\begin{quote}
2026-06-27: Reconstructed proof replacing the original weak version.
\end{quote}

\subsection{1 Overview}\label{overview}

The original Proposition 1 (global expert ranking insufficiency) established the trivial
fact that no single expert is simultaneously optimal across all states
when risk functions cross. This observation is necessary for SCX but not
sufficient as a theoretical foundation. We replace it with a
\textbf{regret lower bound} theorem: any fixed global expert incurs a
quantifiable penalty relative to the state-conditional optimal choice,
and this penalty is bounded from below by a linear function of the
number of states and the minimum risk gap between competing experts.

\begin{center}\rule{0.5\linewidth}{0.5pt}\end{center}

\subsection{2 Notation and Setup}\label{notation-and-setup}

\begin{longtable}[]{@{}
  >{\raggedright\arraybackslash}p{(\linewidth - 2\tabcolsep) * \real{0.4706}}
  >{\raggedright\arraybackslash}p{(\linewidth - 2\tabcolsep) * \real{0.5294}}@{}}
\toprule\noalign{}
\begin{minipage}[b]{\linewidth}\raggedright
Symbol
\end{minipage} & \begin{minipage}[b]{\linewidth}\raggedright
Meaning
\end{minipage} \\
\midrule\noalign{}
\endhead
\bottomrule\noalign{}
\endlastfoot
\(\mathcal{S} = \{s_1,\dots,s_K\}\) & State partition of the input
space \\
\(P(s) > 0\) & Probability mass of state \(s\),
\(\sum_{s \in \mathcal{S}} P(s) = 1\) \\
\(\mathcal{M} = \{1,\dots,M\}\) & Expert indices \\
\(R_m(s)\) & Expected risk of expert \(m\) in state \(s\) \\
\(\ell: \mathcal{Y} \times \mathcal{Y} \to \mathbb{R}^+\) & Loss
function, assumed bounded: \(\ell \in [0, L_{\max}]\) \\
\end{longtable}

\textbf{Assumption 1 (State partition).} The state space \(\mathcal{S}\)
is a finite measurable partition of the input space \(\mathcal{X}\)
induced by a map \(\Pi: \mathcal{X} \to \mathcal{S}\). Within each state
\(s\), the conditional distribution \(P_{X|s}\) governs the input-output
relationship.

\textbf{Assumption 2 (Risk identifiability).} For each state \(s\) and
each expert \(m\), the risk
\(R_m(s) = \mathbb{E}_{x \sim P_{X|s}}[\ell(f_m(x), f^*(x))]\) exists
and is finite. Without loss of generality, all risks are bounded to
\([0, L_{\max}]\).

\begin{center}\rule{0.5\linewidth}{0.5pt}\end{center}

\subsection{3 Definitions}\label{definitions}

\textbf{Definition 1 (Per-state expert ordering).} For state \(s\),
order the experts by increasing risk:

\[R_{(1)}(s) \leq R_{(2)}(s) \leq \cdots \leq R_{(M)}(s)\]

where \(R_{(1)}(s) = \min_{m} R_m(s)\) is the best achievable risk in
state \(s\) and \(R_{(2)}(s)\) is the second-best.

\textbf{Definition 2 (State gap).} The gap in state \(s\) is:

\[\delta(s) = R_{(2)}(s) - R_{(1)}(s) \geq 0\]

The minimum gap across all states is:

\[\delta_{\min} = \min_{s \in \mathcal{S}} \delta(s)\]

\textbf{Definition 3 (Global expert regret).} For a fixed global expert
\(m^*\), its regret relative to the state-conditional optimal choice is:

\[\text{Regret}(m^*) = \sum_{s \in \mathcal{S}} P(s) \cdot \bigl[R_{m^*}(s) - \min_{m \in \mathcal{M}} R_m(s)\bigr]\]

This measures the excess risk incurred by committing to a single expert
across all states, compared to the oracle that selects the best expert
per state.

\begin{center}\rule{0.5\linewidth}{0.5pt}\end{center}

\subsection{4 Theorem 1: Regret Lower
Bound}\label{theorem-1-regret-lower-bound}

\subsubsection{4.1 Statement}\label{statement}

Let \(\mathcal{S}\) be a state partition with \(K\) states, and let
\(\delta_{\min}\) be the minimum gap as defined above. For any fixed
global expert \(m^* \in \mathcal{M}\), define its suboptimal state set:

\[\mathcal{S}_{\text{bad}}(m^*) = \bigl\{s \in \mathcal{S} : m^* \notin \arg\min_{m} R_m(s)\bigr\}\]

Then:

\[\text{Regret}(m^*) \geq \sum_{s \in \mathcal{S}_{\text{bad}}(m^*)} P(s) \cdot \delta(s) \geq P(\mathcal{S}_{\text{bad}}(m^*)) \cdot \delta_{\min}\]

where
\(P(\mathcal{S}_{\text{bad}}(m^*)) = \sum_{s \in \mathcal{S}_{\text{bad}}(m^*)} P(s)\)
is the total probability mass of the states where \(m^*\) is suboptimal.

\subsubsection{4.2 Proof}\label{proof}

\textbf{Step 1.} Fix \(m^*\) and consider any state \(s\). If
\(s \in \mathcal{S}_{\text{bad}}(m^*)\), then \(m^*\) is not the optimal
expert in state \(s\). Therefore:

\[R_{m^*}(s) \geq R_{(2)}(s)\]

since the best possible risk in state \(s\) is \(R_{(1)}(s)\) and the
next best is \(R_{(2)}(s)\), and \(m^*\) is worse than at least the best
expert.

\textbf{Step 2.} The per-state regret contribution for
\(s \in \mathcal{S}_{\text{bad}}(m^*)\) is:

\[R_{m^*}(s) - \min_{m} R_m(s) = R_{m^*}(s) - R_{(1)}(s) \geq R_{(2)}(s) - R_{(1)}(s) = \delta(s)\]

The inequality follows from Step 1: \(R_{m^*}(s) \geq R_{(2)}(s)\), and
\(\min_m R_m(s) = R_{(1)}(s)\).

\textbf{Step 3.} For \(s \notin \mathcal{S}_{\text{bad}}(m^*)\), the
regret contribution is zero by definition since \(m^*\) is optimal in
that state. Summing over all states:

\[\begin{aligned}
\text{Regret}(m^*) &= \sum_{s \in \mathcal{S}} P(s) \cdot \bigl[R_{m^*}(s) - \min_m R_m(s)\bigr] \\
&= \sum_{s \in \mathcal{S}_{\text{bad}}(m^*)} P(s) \cdot \bigl[R_{m^*}(s) - \min_m R_m(s)\bigr] \quad (\text{zero contribution elsewhere}) \\
&\geq \sum_{s \in \mathcal{S}_{\text{bad}}(m^*)} P(s) \cdot \delta(s) \quad (\text{by Step 2}) \\
&\geq \delta_{\min} \cdot \sum_{s \in \mathcal{S}_{\text{bad}}(m^*)} P(s) \quad (\text{since } \delta(s) \geq \delta_{\min}) \\
&= P(\mathcal{S}_{\text{bad}}(m^*)) \cdot \delta_{\min}
\end{aligned}\]

This completes the proof. \(\square\)

\subsubsection{4.3 Tightness}\label{tightness}

The bound is tight. Consider a two-state, two-expert system with
\(P(s_1) = P(s_2) = 0.5\), \(R_1(s_1) = 0\), \(R_2(s_1) = \delta\),
\(R_1(s_2) = \delta\), \(R_2(s_2) = 0\). Then
\(\delta_{\min} = \delta\), and for \(m^* = 1\):

\[\text{Regret}(1) = P(s_2) \cdot \delta = 0.5\delta = P(\mathcal{S}_{\text{bad}}(1)) \cdot \delta_{\min}\]

The bound is attained exactly. For \(m^* = 2\), symmetry gives the same
value.

\begin{center}\rule{0.5\linewidth}{0.5pt}\end{center}

\subsection{5 Corollary: Risk Crossing Implies Positive
Regret}\label{corollary-risk-crossing-implies-positive-regret}

\subsubsection{5.1 Statement}\label{statement-1}

If there exist two states \(s_1, s_2 \in \mathcal{S}\) and two experts
\(a, b \in \mathcal{M}\) such that:

\[R_a(s_1) < R_b(s_1) \quad \text{and} \quad R_a(s_2) > R_b(s_2)\]

(i.e., the expert risk functions \textbf{cross}), then for any global
expert \(m^*\):

\[\text{Regret}(m^*) \geq \delta_{\min} \cdot \min\{P(s_1), P(s_2)\} > 0\]

In particular, when the crossing involves the optimal experts (i.e.,
\(a\) is optimal in \(s_1\) and \(b\) is optimal in \(s_2\)), no single
expert can achieve zero regret.

\subsubsection{5.2 Proof}\label{proof-1}

Under the crossing condition, in state \(s_1\) the best expert is at
least as good as \(a\) (since \(R_a(s_1) < R_b(s_1)\)), and in state
\(s_2\) the best expert is at least as good as \(b\). No single expert
\(m^*\) can be simultaneously equal to both \(a\) and \(b\) (since
\(a \neq b\) by the crossing condition), so at least one of \(s_1\) or
\(s_2\) is in \(\mathcal{S}_{\text{bad}}(m^*)\). Therefore
\(P(\mathcal{S}_{\text{bad}}(m^*)) \geq \min\{P(s_1), P(s_2)\}\), and
applying Theorem 1 gives:

\[\text{Regret}(m^*) \geq \delta_{\min} \cdot \min\{P(s_1), P(s_2)\} > 0\]

The strict positivity follows because \(P(s_1), P(s_2) > 0\) (states are
assumed to have positive probability) and \(\delta_{\min} > 0\) when
crossing involves the optimal experts (since \(\delta(s) > 0\) if the
optimal expert in state \(s\) is strictly better than all others).
\(\square\)

\begin{center}\rule{0.5\linewidth}{0.5pt}\end{center}

\subsection{6 Corollary: Co-monotonicity and Zero
Regret}\label{corollary-co-monotonicity-and-zero-regret}

\subsubsection{6.1 Definition
(Co-monotonicity)}\label{definition-co-monotonicity}

A family of risk functions \(\{R_m\}_{m=1}^M\) is \textbf{co-monotonic}
across states if for any two states \(s_1, s_2 \in \mathcal{S}\) and any
two experts \(a, b \in \mathcal{M}\):

\[\bigl(R_a(s_1) - R_b(s_1)\bigr) \cdot \bigl(R_a(s_2) - R_b(s_2)\bigr) \geq 0\]

Equivalently, the ordering of experts by risk is invariant across
states: there exists a total order \(\preceq\) on \(\mathcal{M}\) such
that \(R_a(s) \leq R_b(s)\) for all \(s \in \mathcal{S}\) whenever
\(a \preceq b\).

\subsubsection{6.2 Statement}\label{statement-2}

Let \(m_{\text{opt}}\) be the best global expert:

\[m_{\text{opt}} = \arg\min_{m} \text{Regret}(m)\]

Then \(\text{Regret}(m_{\text{opt}}) = 0\) \textbf{if and only if} there
exists an expert \(m^*\) that is simultaneously optimal in every state:

\[\forall s \in \mathcal{S} : R_{m^*}(s) = \min_{m} R_m(s)\]

A \textbf{sufficient condition} for this is pairwise co-monotonicity of
all expert risk functions: when experts are co-monotonic, the globally
minimum-risk expert under the total order is uniformly optimal across
all states, giving zero regret.

\subsubsection{6.3 Proof}\label{proof-2}

\textbf{(\(\Leftarrow\)).} If an expert \(m^*\) is optimal in every
state, then \(\mathcal{S}_{\text{bad}}(m^*) = \varnothing\) and Theorem
1 gives \(\text{Regret}(m^*) \geq 0\). Direct calculation gives
\(\text{Regret}(m^*) = 0\) since each term
\(R_{m^*}(s) - \min_m R_m(s) = 0\).

\textbf{(\(\Rightarrow\)).} Suppose
\(\text{Regret}(m_{\text{opt}}) = 0\). Then for every state \(s\),
\(P(s) \cdot [R_{m_{\text{opt}}}(s) - \min_m R_m(s)] = 0\). Since
\(P(s) > 0\), we must have \(R_{m_{\text{opt}}}(s) = \min_m R_m(s)\) for
all \(s\), establishing that \(m_{\text{opt}}\) is simultaneously
optimal everywhere.

\textbf{Co-monotonicity as sufficient condition.} Under co-monotonicity,
the experts are totally ordered by risk. Let \(m^*\) be the minimal
expert under this order (i.e., \(R_{m^*}(s) \leq R_m(s)\) for all \(m\)
and all \(s\)). Then \(m^*\) is optimal in every state, and by the
(\(\Leftarrow\)) direction, \(\text{Regret}(m^*) = 0\). \(\square\)

\subsubsection{6.4 Testable Condition for
Co-monotonicity}\label{testable-condition-for-co-monotonicity}

Given empirical risk estimates \(\hat{R}_m(s)\) computed from data,
co-monotonicity can be tested pairwise:

\begin{enumerate}
\def\labelenumi{\arabic{enumi}.}
\item
  \textbf{Pairwise sign test.} For each pair \((a, b)\), compute the
  sign vector:
  \[\sigma_{a,b}(s) = \operatorname{sgn}\bigl(\hat{R}_a(s) - \hat{R}_b(s)\bigr)\]
  for all \(s \in \mathcal{S}\). If for every pair \(\sigma_{a,b}(s)\)
  is constant across \(s\) (all entries are \(+1\), all \(-1\), or all
  \(0\)), then the data are consistent with co-monotonicity.
\item
  \textbf{Rank correlation test.} For each pair \((a, b)\), compute
  Kendall's \(\tau\) or Spearman's \(\rho\) of the risk vectors
  \((\hat{R}_a(s))_{s \in \mathcal{S}}\) and
  \((\hat{R}_b(s))_{s \in \mathcal{S}}\). Under co-monotonicity, all
  pairwise rank correlations equal \(+1\). Deviations indicate risk
  crossing.
\item
  \textbf{Global crossing count.} The number of crossing pairs:
  \[C = \bigl|\{(a, b, s_1, s_2) : (R_a(s_1) - R_b(s_1))(R_a(s_2) - R_b(s_2)) < 0\}\bigr|\]
  quantifies the departure from co-monotonicity. When \(C > 0\), regret
  must be positive for any global expert, and the lower bound of Theorem
  1 applies.
\end{enumerate}

\begin{center}\rule{0.5\linewidth}{0.5pt}\end{center}

\subsection{7 Empirical Regret Bound}\label{empirical-regret-bound}

In practice, risks are estimated from finite data. Let \(\hat{R}_m(s)\)
be the empirical risk estimate based on \(n_s\) samples in state \(s\).

\subsubsection{7.1 Statement}\label{statement-3}

With probability at least \(1 - \eta\), for every expert \(m\), the
regret satisfies:

\[\text{Regret}(m) \geq \sum_{s \in \mathcal{S}} P(s) \cdot \bigl[\hat{R}_m(s) - \min_{m'} \hat{R}_{m'}(s)\bigr] - L_{\max} \sum_{s \in \mathcal{S}} P(s) \cdot \sqrt{\frac{2\log(2M/\eta)}{n_s}}\]

\subsubsection{7.2 Proof Sketch}\label{proof-sketch}

By the Hoeffding inequality, for each state \(s\) and expert \(m\):

\[|\hat{R}_m(s) - R_m(s)| \leq L_{\max} \sqrt{\frac{2\log(2M/\eta)}{n_s}}\]

with probability \(1 - \eta/M\). A union bound over \(M\) experts and
\(K\) states, combined with the decomposition of regret, yields the
result. This bound allows practitioners to compute a
\textbf{high-probability lower bound} on regret from finite-sample risk
estimates. \(\square\)

\begin{center}\rule{0.5\linewidth}{0.5pt}\end{center}

\subsection{8 Implications for SCX}\label{implications-for-scx}

\begin{enumerate}
\def\labelenumi{\arabic{enumi}.}
\item
  \textbf{Regret as a design principle.} Theorem 1 quantifies the
  penalty for ignoring state structure. The lower bound
  \(P(\mathcal{S}_{\text{bad}}(m^*)) \cdot \delta_{\min}\) shows that
  regret grows linearly with both the mass of suboptimal states and the
  minimum risk gap.
\item
  \textbf{State discovery drives regret reduction.} SCX's state
  discovery directly reduces \(P(\mathcal{S}_{\text{bad}}(m^*))\) by
  identifying regions where different experts excel. A perfect state
  partition (where each state has a single optimal expert) drives regret
  to zero for the state-conditioned system.
\item
  \textbf{Co-monotonicity as a diagnostic.} The co-monotonicity test of
  Section 6.4 provides a data-driven criterion for whether a global
  expert is sufficient. When the test fails, SCX's state-conditioned
  approach is guaranteed to improve over any fixed expert.
\item
  \textbf{Sample-efficient testing.} The empirical bound of Section 7
  enables practitioners to compute statistically valid lower bounds on
  regret with finite data, guarding against false conclusions due to
  estimation noise.
\end{enumerate}

\begin{center}\rule{0.5\linewidth}{0.5pt}\end{center}

\subsection{References}\label{references}

\begin{enumerate}
\def\labelenumi{\arabic{enumi}.}
\tightlist
\item
  SCX theoretical framework. \texttt{01\_SCX\_Core\_Framework\_Math\_Analysis.md}
  Section 2.1
\item
  SCX mathematical foundations. \texttt{05\_Math\_Roots\_and\_Proofs.md}
  Proposition 1
\item
  Tsybakov, A. B. (2004). Optimal aggregation of classifiers in
  statistical learning. \emph{Annals of Statistics}, 32(1), 135-166.
  {[}Related: minimax lower bounds for aggregation{]}
\item
  Hoeffding, W. (1963). Probability inequalities for sums of bounded
  random variables. \emph{Journal of the American Statistical
  Association}.
\end{enumerate}

\end{document}
